\chapter{Let the Test Own Its Product}

The black-box test currently depends on the FreeBSD Handbook record that we
created by hand. That is useful while learning the path through the system,
but it is weak test isolation: deleting or renaming that old record would make
an otherwise-correct deployment fail.

We will make the test create a unique product through the real catalog API,
observe it through nginx and store, and delete it afterward. The fixture will
therefore use the same database and HTTP boundaries as the application rather
than reaching directly into PostgreSQL.

\section{Exercise 25A: Replace the static fixture}

Open a root console in the \texttt{test} jail and edit the existing test:

\begin{lstlisting}[language=sh]
vi /usr/local/tests/store_e2e_test.rb
\end{lstlisting}

Replace its contents with:

\begin{lstlisting}[language=Ruby]
require "json"
require "minitest/autorun"
require "net/http"
require "uri"

class StoreEndToEndTest < Minitest::Test
  STORE_URL = ENV.fetch(
    "STORE_URL", "http://10.17.89.26/"
  )
  CATALOG_URL = ENV.fetch(
    "CATALOG_URL", "http://10.17.89.24:3000"
  )

  def setup
    token = [
      Process.pid, Time.now.to_i, rand(1_000_000)
    ].join("-")
    @sku = "E2E-#{token}"
    @name = "End-to-End Product #{token}"

    uri = URI.join("#{CATALOG_URL}/", "products")
    request = Net::HTTP::Post.new(uri)
    request["Content-Type"] = "application/json"
    request.body = JSON.generate(
      product: {
        sku: @sku,
        name: @name,
        price_cents: 12_345
      }
    )

    response = perform(uri, request)
    assert_equal "201", response.code, response.body
    @product_id = JSON.parse(response.body).fetch("id")
  end

  def teardown
    return unless @product_id

    uri = URI.join(
      "#{CATALOG_URL}/", "products/#{@product_id}"
    )
    response = perform(uri, Net::HTTP::Delete.new(uri))
    assert_equal "204", response.code, response.body
  end

  def test_product_is_visible_through_the_web_edge
    uri = URI(STORE_URL)
    response = perform(uri, Net::HTTP::Get.new(uri))

    assert_equal "200", response.code, response.body
    assert_includes response.body, @name
    assert_includes response.body, @sku
    assert_includes response.body, "$123.45"
  end

  private

  def perform(uri, request)
    Net::HTTP.start(
      uri.host,
      uri.port,
      open_timeout: 2,
      read_timeout: 5
    ) do |http|
      http.request(request)
    end
  end
end
\end{lstlisting}

The process ID, current time, and random suffix make the name and SKU distinct
for each invocation. The exact values do not matter; the test remembers them
and looks for those values in the final HTML.

The catalog request belongs in \texttt{setup} because the product must exist
before store renders the page. Minitest runs \texttt{teardown} after the test
even when an assertion fails, so the corresponding API deletion is the best
place for normal cleanup. If creation fails before an ID is returned, the
guard prevents a meaningless delete request.

\section{Exercise 25B: Check the Ruby syntax}

Inside the test jail, ask Ruby to parse the file without running it:

\begin{lstlisting}[language=sh]
ruby -c /usr/local/tests/store_e2e_test.rb
\end{lstlisting}

Expect \texttt{Syntax OK}. This catches copying mistakes without creating a
catalog product or making a request through nginx.

\section{Exercise 25C: Run the owned-fixture test}

Leave the jail console and invoke the host runner:

\begin{lstlisting}[language=sh]
mdo /usr/local/sbin/rails-on-jails-test
\end{lstlisting}

The runner starts any missing jails before executing the updated file. Expect
one run, nine assertions, no failures, and no errors. Those assertions cover
catalog creation, the rendered response, and catalog deletion.

Run the same command a second time. It should pass with a different generated
name and SKU. This demonstrates that the test does not depend on a record left
by its previous run.

\section{Checkpoint}

Stop after two consecutive host-runner invocations pass. The test now owns the
data that determines its result and cleans up through the public catalog API.
The original FreeBSD Handbook can remain as demonstration data, but it is no
longer part of the automated test contract.
